\section*{Análisis Bibliográfico Detallado}

Este documento presenta un análisis detallado de todas las fuentes bibliográficas utilizadas en la propuesta de tesis, justificando su selección, relevancia y aplicación en el documento.

\subsection*{Akenine-Möller, T., Haines, E., \& Hoffman, N. (2018)}
\textbf{Cita:} Akenine-Möller, T., Haines, E., \& Hoffman, N. (2018). \textit{Real-Time Rendering} (4th ed.). CRC Press. \\
\textbf{Año:} 2018 \\
\textbf{Resumen:} La referencia estándar en la industria para algoritmos de renderizado en tiempo real. \\
\textbf{Relevancia:} Base técnica para el pipeline gráfico. \\
\textbf{Uso:} Marco Teórico (Fundamentos de Renderizado).

\subsection*{Bartlett, K. A., \& Dorribo Camba, J. (2023)}
\textbf{Cita:} Bartlett, K. A., \& Dorribo Camba, J. (2023). The Role of a Graphical Interpretation Factor in the Assessment of Spatial Visualization. \textit{Spatial Cognition \& Computation}. \\
\textbf{Año:} 2023 \\
\textbf{Resumen:} Estudio sobre cómo la interpretación gráfica afecta la evaluación espacial. \\
\textbf{Relevancia:} Evidencia reciente de la importancia de visualizaciones claras. \\
\textbf{Uso:} Marco Teórico (Visualización Espacial).

\subsection*{Bowman, D. A., et al. (2004)}
\textbf{Cita:} Bowman, D. A., et al. (2004). \textit{3D User Interfaces: Theory and Practice}. Addison-Wesley. \\
\textbf{Año:} 2004 \\
\textbf{Resumen:} Teoría y práctica de interfaces 3D. \\
\textbf{Relevancia:} Guía el diseño de la interacción del prototipo. \\
\textbf{Uso:} Marco Teórico y Metodología.

\subsection*{Burley, B. (2012)}
\textbf{Cita:} Burley, B. (2012). Physically-Based Shading at Disney. \textit{SIGGRAPH 2012}. \\
\textbf{Año:} 2012 \\
\textbf{Resumen:} Introducción del modelo PBR de Disney. \\
\textbf{Relevancia:} Base del modelo de iluminación usado en Unity URP. \\
\textbf{Uso:} Marco Teórico (PBR).

\subsection*{Cook, R. L., \& Torrance, K. E. (1982)}
\textbf{Cita:} Cook, R. L., \& Torrance, K. E. (1982). A reflectance model for computer graphics. \textit{ACM TOG}. \\
\textbf{Año:} 1982 \\
\textbf{Resumen:} Modelo de reflectancia fundacional para gráficos por computadora. \\
\textbf{Relevancia:} Fundamento físico del renderizado. \\
\textbf{Uso:} Marco Teórico.

\subsection*{Cowan, N. (2001)}
\textbf{Cita:} Cowan, N. (2001). The magical number 4 in short-term memory. \textit{Behavioral and Brain Sciences}. \\
\textbf{Año:} 2001 \\
\textbf{Resumen:} Revisión de la capacidad de memoria de trabajo (4 items). \\
\textbf{Relevancia:} Justifica la simplificación de la interfaz. \\
\textbf{Uso:} Marco Teórico (Carga Cognitiva).

\subsection*{Darken, R. P., \& Sibert, J. L. (1996)}
\textbf{Cita:} Darken, R. P., \& Sibert, J. L. (1996). Wayfinding Strategies and Behaviours in Large Virtual Worlds. \textit{CHI '96}. \\
\textbf{Año:} 1996 \\
\textbf{Resumen:} Estrategias de navegación en mundos virtuales. \\
\textbf{Relevancia:} Diseño de navegación para evitar desorientación. \\
\textbf{Uso:} Marco Teórico.

\subsection*{Fransson, E., et al. (2024)}
\textbf{Cita:} Fransson, E., et al. (2024). A Comparison of Performance on WebGPU and WebGL. \textit{IEEE CoG 2024}. \\
\textbf{Año:} 2024 \\
\textbf{Resumen:} Comparativa de rendimiento WebGL vs WebGPU. \\
\textbf{Relevancia:} Justificación tecnológica y visión a futuro. \\
\textbf{Uso:} Estado del Arte.

\subsection*{Hegarty, M. (2004)}
\textbf{Cita:} Hegarty, M. (2004). Mechanical reasoning by mental simulation. \textit{Trends in Cognitive Sciences}. \\
\textbf{Año:} 2004 \\
\textbf{Resumen:} Razonamiento mecánico y simulación mental. \\
\textbf{Relevancia:} Justifica el uso de animación para entender mecanismos. \\
\textbf{Uso:} Marco Teórico.

\subsection*{Hegarty, M., \& Waller, D. (2004)}
\textbf{Cita:} Hegarty, M., \& Waller, D. (2004). Spatial Abilities at Different Scales. \textit{Intelligence}. \\
\textbf{Año:} 2004 \\
\textbf{Resumen:} Diferencias individuales en habilidades espaciales. \\
\textbf{Relevancia:} Consideración de diferencias de usuario en el diseño. \\
\textbf{Uso:} Marco Teórico.

\subsection*{Hevner, A. R., et al. (2004)}
\textbf{Cita:} Hevner, A. R., et al. (2004). Design Science in Information Systems Research. \textit{MIS Quarterly}. \\
\textbf{Año:} 2004 \\
\textbf{Resumen:} Marco metodológico DSR. \\
\textbf{Relevancia:} Metodología base del proyecto. \\
\textbf{Uso:} Metodología.

\subsection*{Hutchins, E. (1995)}
\textbf{Cita:} Hutchins, E. (1995). \textit{Cognition in the Wild}. MIT Press. \\
\textbf{Año:} 1995 \\
\textbf{Resumen:} Cognición distribuida en entornos reales. \\
\textbf{Relevancia:} Contexto de cómo las herramientas afectan la cognición. \\
\textbf{Uso:} Marco Conceptual.

\subsection*{Karis, B. (2013)}
\textbf{Cita:} Karis, B. (2013). Real Shading in Unreal Engine 4. \textit{SIGGRAPH 2013}. \\
\textbf{Año:} 2013 \\
\textbf{Resumen:} Adaptación del modelo Disney para tiempo real en juegos. \\
\textbf{Relevancia:} Referencia técnica para aproximaciones PBR en tiempo real. \\
\textbf{Uso:} Marco Teórico.

\subsection*{Khronos Group (2011)}
\textbf{Cita:} Khronos Group. (2011). \textit{WebGL Specification 1.0}. \\
\textbf{Año:} 2011 \\
\textbf{Resumen:} Especificación estándar de WebGL. \\
\textbf{Relevancia:} Definición tecnológica base. \\
\textbf{Uso:} Estado del Arte.

\subsection*{Köhler, W. (1929)}
\textbf{Cita:} Köhler, W. (1929). \textit{Gestalt Psychology}. Liveright. \\
\textbf{Año:} 1929 \\
\textbf{Resumen:} Principios de la psicología de la Gestalt. \\
\textbf{Relevancia:} Fundamento de percepción visual (agrupamiento, cierre). \\
\textbf{Uso:} Marco Teórico (Percepción).

\subsection*{Mayer, R. E. (2005, 2021)}
\textbf{Cita:} Mayer, R. E. (2005/2021). \textit{The Cambridge Handbook of Multimedia Learning}. \\
\textbf{Año:} 2005, 2021 \\
\textbf{Resumen:} Teoría Cognitiva del Aprendizaje Multimedia. \\
\textbf{Relevancia:} Principios de diseño instruccional (contigüidad, coherencia). \\
\textbf{Uso:} Marco Teórico y Metodología.

\subsection*{Miller, G. A. (1956)}
\textbf{Cita:} Miller, G. A. (1956). The Magical Number Seven, Plus or Minus Two. \textit{Psychological Review}. \\
\textbf{Año:} 1956 \\
\textbf{Resumen:} Capacidad limitada de procesamiento de información. \\
\textbf{Relevancia:} Antecedente clásico de la carga cognitiva. \\
\textbf{Uso:} Marco Teórico.

\subsection*{Cabello, R. [mrdoob] (2024)}
\textbf{Cita:} Cabello, R. (2024). \textit{three.js - JavaScript 3D Library}. \\
\textbf{Año:} 2024 \\
\textbf{Resumen:} Biblioteca Three.js. \\
\textbf{Relevancia:} Benchmarking tecnológico. \\
\textbf{Uso:} Estado del Arte.

\subsection*{Nielsen, J. (1994, 2000)}
\textbf{Cita:} Nielsen, J. (1994). \textit{Usability Engineering}; (2000). Why You Only Need to Test with 5 Users. \\
\textbf{Año:} 1994, 2000 \\
\textbf{Resumen:} Heurísticas de usabilidad y tamaño de muestra. \\
\textbf{Relevancia:} Metodología de validación. \\
\textbf{Uso:} Metodología.

\subsection*{Norman, D. A. (2013)}
\textbf{Cita:} Norman, D. A. (2013). \textit{The Design of Everyday Things}. \\
\textbf{Año:} 2013 \\
\textbf{Resumen:} Diseño centrado en el usuario. \\
\textbf{Relevancia:} Principios de interacción. \\
\textbf{Uso:} Marco Teórico.

\subsection*{Paivio, A. (1986)}
\textbf{Cita:} Paivio, A. (1986). \textit{Mental Representations: A Dual Coding Approach}. \\
\textbf{Año:} 1986 \\
\textbf{Resumen:} Teoría de Codificación Dual. \\
\textbf{Relevancia:} Justificación de uso de texto + imagen 3D. \\
\textbf{Uso:} Marco Teórico.

\subsection*{Ries, E. (2011)}
\textbf{Cita:} Ries, E. (2011). \textit{The Lean Startup}. \\
\textbf{Año:} 2011 \\
\textbf{Resumen:} Metodología Lean y MVP. \\
\textbf{Relevancia:} Estrategia de desarrollo del prototipo. \\
\textbf{Uso:} Metodología.

\subsection*{Schlick, C. (1994)}
\textbf{Cita:} Schlick, C. (1994). An Inexpensive BRDF Model for Physically-based Rendering. \textit{Computer Graphics Forum}. \\
\textbf{Año:} 1994 \\
\textbf{Resumen:} Aproximación eficiente para cálculos de reflectancia (Fresnel). \\
\textbf{Relevancia:} Técnica de optimización usada en shaders PBR en tiempo real. \\
\textbf{Uso:} Marco Teórico.

\subsection*{Sweller, J. (1988, 2019)}
\textbf{Cita:} Sweller, J. (1988, 2019). Cognitive load theory references. \\
\textbf{Año:} 1988, 2019 \\
\textbf{Resumen:} Teoría de Carga Cognitiva. \\
\textbf{Relevancia:} Fundamento teórico principal del problema de investigación. \\
\textbf{Uso:} Marco Teórico.

\subsection*{Unity Technologies (2021)}
\textbf{Cita:} Unity Technologies. (2021). \textit{Advances in Real-Time Rendering in Games}. SIGGRAPH Course. \\
\textbf{Año:} 2021 \\
\textbf{Resumen:} Avances recientes en renderizado de Unity. \\
\textbf{Relevancia:} Estado del arte en motores de juego. \\
\textbf{Uso:} Estado del Arte.

\subsection*{Unity Technologies (2024)}
\textbf{Cita:} Documentación oficial (WebGL, Memory). \\
\textbf{Año:} 2024 \\
\textbf{Resumen:} Manuales técnicos de Unity. \\
\textbf{Relevancia:} Referencia técnica de implementación. \\
\textbf{Uso:} Estado del Arte y Metodología.

\subsection*{Walter, B., et al. (2007)}
\textbf{Cita:} Walter, B., et al. (2007). Microfacet Models for Refraction through Rough Surfaces. \textit{EGSR}. \\
\textbf{Año:} 2007 \\
\textbf{Resumen:} Modelos de microfacetas para superficies rugosas (GGX). \\
\textbf{Relevancia:} Base de la función de distribución normal (NDF) usada en PBR moderno. \\
\textbf{Uso:} Marco Teórico.

\subsection*{Ware, C. (2012)}
\textbf{Cita:} Ware, C. (2012). \textit{Information Visualization: Perception for Design}. Morgan Kaufmann. \\
\textbf{Año:} 2012 \\
\textbf{Resumen:} Percepción visual aplicada al diseño de información. \\
\textbf{Relevancia:} Guía para la presentación efectiva de datos visuales. \\
\textbf{Uso:} Marco Teórico.

\subsection*{Wertheimer, M. (1923)}
\textbf{Cita:} Wertheimer, M. (1923). Untersuchungen zur Lehre von der Gestalt II. \\
\textbf{Año:} 1923 \\
\textbf{Resumen:} Obra fundacional de la Gestalt. \\
\textbf{Relevancia:} Principios perceptuales básicos. \\
\textbf{Uso:} Marco Teórico.

\subsection*{Yu, G., et al. (2023)}
\textbf{Cita:} Yu, G., et al. (2023). A survey of real-time rendering on Web3D application. \\
\textbf{Año:} 2023 \\
\textbf{Resumen:} Survey reciente de Web3D. \\
\textbf{Relevancia:} Contextualización del estado del arte. \\
\textbf{Uso:} Estado del Arte.
